% !TEX root = ../main.tex
\section{Realization Details on Cloud AI 100}
\label{app:impl}

This appendix records how the Cloud AI 100 export and compiler shaped the realization summarized in \S\ref{sec:impl}.

\subsection{What the Compiler Does by Default}
\label{app:impl-default}

\textbf{Slicing across SoCs.} With its default heuristic the partitioner splits the down-projection weights of all 64
lanes four ways along the output dimension, so each SoC computes a 512-column slice of every expert's output. This
requires an all-gather of the intermediate activation (36\,MiB per stage at $C=128$) and a re-shard of the output
(24\,MiB per stage) before the combine; cross-SoC traffic reaches 121.6\,MiB per layer, while gate and up projections stay
expert-parallel.

\textbf{A dense accumulator.} The exported combine scatters each stage's weighted expert rows into an accumulator of
shape $[64, T, d]$, zero-filled by a single core, gathered and updated between the two stages, reduced over each SoC's
16 lanes by one einsum that the compiler runs through DRAM on one core, and finally summed over the four SoCs on four
cores of SoC~0. From $T=256$ on, this last sum also spills into a DRAM-backed reduction on one core (0.84\,ms). Every
part scales with $T$; none scales with the number of routed rows.

\textbf{Lane counts must divide 16.} The compiler maps a stage onto a SoC's cores by its lanes per SoC: counts that divide
16 map evenly, while in our lane-split programs 17--19 lanes split the GEMMs unevenly and 20--32 lanes collapsed onto four
cores per SoC (3.7--4.3\,ms for a stage that runs in 0.5\,ms on 16 cores).

\textbf{Stages run one after another.} The compiler schedules a layer's expert stages sequentially, in its own order;
in our tiered programs a stage of $L$ lanes per SoC ran one lane per core on $L$ cores. Every additional stage adds a
window of its own.

\subsection{Expert-Parallel Compilation with an On-Device KV Cache}
\label{sec:ep}

We compile with a weight-splitting degree of one across SoCs (\texttt{-mdts-mos=1}), which forbids cutting any weight
matrix across SoCs. The lane axis is then the only way to distribute the batched expert GEMMs, and each SoC computes 16
whole experts per stage, one per core; the all-gather and re-shard of the down projection disappear, and the layer's
cross-SoC traffic falls from 121.6 to 1.5\,MiB.

The setting is global and also stops the compiler from splitting the attention projections by KV head. The cache update
then comes out token-sliced while its input is whole, the compiler cannot pair the 96 cache inputs with their outputs,
and all 192 KV buffers become host I/O. We restore the pairing at the graph level with \emph{head-parallel attention}:
the four KV-head groups of eight query heads each \citep{ainslie2023gqa} become a leading batch axis of every attention
operator, with per-group weights of shapes $[4,2048,1024]$ for queries, $[4,2048,128]$ for keys and values and
$[4,1024,2048]$ for the output projection, whose four partial outputs are summed at the end. Under expert-parallel
compilation the batch axis lands one group per SoC, the cache operations of KV head $h$ stay on SoC $h$, and the program
again exposes no KV buffer; the regrouped output projection changes only the FP16 summation order.

\textbf{The group sum.} The four partial outputs live on four SoCs. Our first version summed them with an einsum over the
group axis, which the compiler ran on few cores: in a full-model profile of \sys at $T=512$ it took 2.39 of each layer's
7.73\,ms (31\%). Written as the elementwise tree $(G_0+G_1)+(G_2+G_3)$ per token tile, like the MoE layer's final sum, it
spreads over all cores and saves 17, 67 and 130\,ms per chunk in the naive program at $T=128$, 256 and 512. This repair
belongs to attention; the evaluation keeps it as step 2$'$ and uses it in every later program.

\subsection{Lowering the Combine}
\label{sec:impl-comb}

\textbf{Dense-path rewrites.} Before the combine is replaced, three rewrites remove serialized work.
(i) The first stage's outputs were added to rows gathered from the freshly zeroed accumulator; we drop this read of
zeros. (ii) The per-SoC sum over lanes, one einsum that the compiler ran through DRAM on one core, is cut into 16 tiles of
$T/16$ tokens. (iii) The routing's prefix sums over tokens, which give each assignment its row in the expert's compact
input and ran as a serial \texttt{CumSum} on one core, become Hillis--Steele scans of $\log_2 T$ shift-and-add steps
\citep{hillis1986data}. All three leave the output bit-identical. They repair this export and this compiler's lowering
rather than change the combine's formulation, and the evaluation keeps them as steps of their own: 3a for the scans, 3b
for (i) and (ii).

\textbf{Ownership-preserving aggregation.} The device gather does not zero-fill out-of-range indices, so a safe index
and a mask are both required, and the per-SoC work must stay on the lane axis the compiler already splits over SoCs:
written with global row indices or as explicit per-SoC slices, the same sum shipped 1539\,MiB of expert outputs between
SoCs per layer and made the full model 3.4--5.3$\times$ slower (\S\ref{sec:mot-own}).

\textbf{Elementwise final sum.} Written as 16 einsum tiles, the sum over the four partials ran on four cores per SoC;
written as the elementwise tree $(P_0+P_1)+(P_2+P_3)$ per token tile, it spreads over all cores, and the compiler lowers it
as a reduce-scatter that the next layer's input normalization completes.

\subsection{Lowering the Rank Sort and the Tiers}
\label{sec:impl-cap}

\textbf{One gathered constant.} The counts come from the routing mask, and the per-SoC sort is a \texttt{TopK} over each
SoC's 32 counts. Each tier's three weight matrices come from \texttt{Gather(bank,$\,\pi$)} on one constant that holds all
128 experts, the only gather form this compiler keeps fast, lowered to indirect DMA (\S\ref{sec:mot-own}). The constant
is replicated, so every SoC holds the weights of all 128 experts, 22.5\,GiB per SoC instead of 6.7\,GiB, within the SoC's
32\,GB; the algorithm itself needs only each SoC's own 32 experts.

\textbf{Off the critical path.} The routing mask, the prefix scans and the slot tables are computed once over all 128
experts in native order, independent of the sort; only a row permutation of these tables and the weight gathers wait for
it. Sorting first and running the routing on permuted columns makes everything wait (12 layers: +2.5\,ms over
capacity-$T$ padding, against +1.6\,ms). Gathered weights are also sensitive to the compiler's maximum tile size: with the
default, the first expert GEMM waits 0.085\,ms for its first gathered tile, with 512\,KiB tiles 0.032\,ms. For the rank sort the best
size is 512\,KiB for $T\le256$ and 1024\,KiB for $T=512$ (Appendix~\ref{app:tiles}); other programs differ (Table~\ref{tab:attr}).

\textbf{Tier widths and count.} Stages map evenly only when their lanes per SoC divide 16, run one after another, and
use $L$ cores for $L$ lanes (Appendix~\ref{app:impl-default}), so tiers take widths from $\{1,2,4,8,16\}$. At $T=512$ the first tier
is bound by the SoC's shared memory traffic: eight lanes of 512 rows finish in 0.8--1.0\,ms, sixteen in 2.2\,ms. We
deploy three tiers of 8, 8 and 16 lanes (\S\ref{sec:calib}, \S\ref{sec:eval-pad}).

\subsection{Graph Transformations}
\label{sec:impl-graph}

\sys is a set of ONNX graph transformations, about 1,200 lines of Python, applied to the Qwen3-30B-A3B graph that the
vendor's QEfficient library exports for Cloud AI 100 \citep{qefficient}. Each step of the ladder (\S\ref{sec:setup}) is a
separate pass, so every intermediate program can be compiled and measured. The passes locate operators by structure rather than
by name (layer 0 numbers some nodes differently and folds a 64-lane constant into its stage), regenerate the
$T$-dependent constants of the scans and tiles for each chunk length, and clone the stage subgraph once per tier with its
own row order, weight gathers and capacity. A static \texttt{If} in the export that squeezes the token table is replaced
by the equivalent \texttt{Squeeze} in every clone, and the aggregation's map from an expert to (tier, SoC, lane) becomes a set of
small per-tier tables. The passes also build the evaluation's variants: the export's own combine, the aggregation formulations of
Table~\ref{tab:own}, fixed expert order
through the gathers, fixed-size blocks, and tier capacities read from input lengths for the overflow fallback.

Each chunk length is its own program, with sequence length $T$ and context $2T$, MXFP6 matrix weights, FP16
activations, 16 cores per SoC, depth-first scheduling, expert-parallel slicing and retained state
(Appendix~\ref{app:flags}). Static programs occupy 25\,GiB, programs with gathered weights 88\,GiB. Every program returns each lane's token count in every layer.
